\scanpage{176}
(cf p.ex. Sém. Chevalley 58/59, exp. 6 p. 06)
\[
 (\Delta_i^0+c_i)^p=\underset{\substack{\Vert\\0\ \text{par hyp}}}{\Delta_i^{0\,p}}+c_i^p+D_i^{p-1}c_i
\]
donc la condition envisagée s'écrit
\[
 c_i^p+D_i^{p-1}c_i=0\qquad i\in I
 \tag*{(*)}
\]
Or écrivons que \inserted{$d\omega=0$,} $C\omega=\widetilde\omega$, \deleted{on trouve} $(\omega=\sum c_i\,dx_i)$ on trouve pour $d\omega=0$ la condition
\[
 D_i c_j-D_j c_i=0\qquad i,j\in I
 \tag*{(**)}
\]
et on peut alors écrire
\[
 \widetilde C\omega=
 \deleted{\sum_i[(c_i^p\unsure{+|-}D_i^{p-1}c_i)\otimes_B1_\Lambda]\,d\widetilde x_i}\ 
 \sum_i-\underbrace{D_i^{p-1}c_i}_{\substack{\text{est dans }\Lambda\\
 \text{car égal à}\\\langle\widetilde D_i,C\omega\rangle1_\Lambda}}
 (d\widetilde x_i\otimes1_\Lambda
 \tag*{(***)}
\]
\deleted{ce qui signifie aussi que $\sum(c_i^p+D_i^{p-1}c_i)x_i^{p-1}dx_i\overset{\text{cohomologue}}{\sim}\sum c_i\,dx_i,$ i.e. que $\sum c_i^p$}
[grâce à la formule de \unsure{transposition|transformation} de loc cit. p. 05 (prop. 3)].
% slanted
\margin{inclure dans le formulaire du \S\,3 ou \S\,4}

alors qu'on a $\widetilde\omega=\sum c_i^p(d\widetilde x_i\otimes1_\Lambda)$

et on trouve donc,
\[
 -\widetilde C\omega+\overline\omega=\sum(c_i^p+D_i^{p-1}c_i)(d\widetilde x_i\otimes1_\Lambda)
 \tag*{$*$}
\]
donc c'est nul ssi (*) l'est, cqfd.

Nous utilisons maintenant le fait bien connu qu'une connexion à courbure nulle \deleted{\unsure{pour|dont}} et compatible à la puissance $p$-ième, sur un Module quasi-cohérent, équivaut à une \uline{donnée de descente} de celui-ci \deleted{\unsure{de}} $X$ à $(X^{p/S},\underline\Lambda)$. \deleted{On trouve que}
% archive sheet number: 175

